%% title: Training: two recipes, one to be kept
%% target: Sections/Methods.tex
%% anchor: \subsection{Loss function}\label{sec:loss_function}
%% action: before
%% why: The two drafts trained with different recipes. Until one is chosen, every result has to say which one it used. The lepinet recipe is the baseline of every shift, self-training and long-tail result.
%% source: configs/20260729_ucloud_singlehead_species_effnetv2s.yaml; src/lepinet/data.py (sample_weights); docs/design-decisions.md section 1 for the ladder
%% check: ON HOLD: rewritten once the decision sheet is pinned (journal/2026-10-23-which-recipe-does-the-merged-paper-rerun.md); the owner gave priority to the colleague's settings. | \cref{sec:results_heads} arrives with a results step; Overleaf shows ?? until then. | Asger's settings for the head comparison (epochs, image size, optimiser): marked X. | The recipe names (M and L) are placeholders.
%% status: hold
\subsection{Training}\label{sec:training}
% TODO(GM, AS): one of the two recipes will be kept for the final paper. Until then every table and figure names the recipe it used.
Two training recipes are used in this draft, and every result names the one it used.

\textbf{Recipe M} (mini\_trainer) trains with the EMLA loss (\cref{sec:loss_function}) and the prototype regulariser (\cref{sec:regularization}) for \textbf{X} epochs at \textbf{X} pixels with \textbf{X}.
It is used for the comparison of head architectures (\cref{sec:results_heads}).

\textbf{Recipe L} (lepinet) trains EfficientNetV2-S for 5 epochs with plain cross-entropy and square-root class oversampling: each image of a species with $n$ training images is drawn with weight $n^{-1/2}$, so the species as a whole carries weight $n^{1/2}$.
The backbone is optimised with Muon and the classification head with AdamW, under a one-cycle schedule with half an epoch of warm-up, a learning rate of $10^{-3}$, batches of 64 and gradient clipping at 5.
Images are resized to 460 pixels and augmented by flips, rotations up to $15^\circ$ and zoom up to 1.1, then cropped to 256 pixels.
Training uses fp16 mixed precision with the head kept in fp32; heads with an additive angular margin overflow in fp16 and are trained in bf16.
Recipe L is used for all results on domain shift, self-training and the long tail.
